\section{Audit findings, corrections, and integration consequences}
\label{sec:audit}

The source formulas used in the proof were rederived at the interfaces
where normalization mistakes are most likely: the omitted-binomial
convention in $M_{a,b}$, the factor two in the Bernoulli coordinates
$U_j=2B_{2b-2j}$, the retained-column reversal, and the scalar
multiplier in \eqref{eq:exact-multiplier}.  No mathematical
contradiction was found in those inherited formulas.  The following
status and presentation corrections should nevertheless be made when
this package is integrated.

\subsection{Research-status updates}

\begin{enumerate}[label=\textbf{A\arabic*.}]
\item The statement that the smaller-order-ten family remains open is
now obsolete.  It should be replaced by
\cref{thm:complete-order-ten}: every $M_{5,b}$ with $b\ge6$ is
nonsingular.

\item The exceptional congruence class
$b\equiv962\pmod{2048}$ in the preceding $2$-adic theorem was a
limitation of that congruence, not a candidate zero set.  The new finite
certificate includes every member of this class below $1920$, beginning
with $b=962$, and the inherited eventual theorem covers all members at
or above $1920$.

\item The broad unequal-order conjecture is still open.  Completion of
$p=10$ must not be described as a proof for all $p<q$.  The next fixed
smaller order is $p=12$, and the all-order problem requires control
uniform in $a$.
\end{enumerate}

\subsection{Certificate wording}

The residue vector is deliberately mixed: for $b<10$ it contains the
direct source determinant, while for $b\ge10$ it contains the normalized
five-square determinant.  Calling every entry ``the residue of
$\det M_{5,b}$'' would be inaccurate.  The exact definition is
\eqref{eq:proof-residue}; equivalence to source nonvanishing follows
from the nonzero multiplier.  The JSON metadata and integration notes
preserve this distinction.

A checksum alone is not a certificate.  The proof objects are the full
residue vector, the formulas generating each matrix, the denominator
argument, and the replay programs.  The SHA-256 value only detects
accidental or silent modification of that data.

\subsection{Analytic normalization}

The ordinary identities require the centered coordinate
$x_n=n+1/2$.  At this center the subtraction polynomial contains only
even powers and every such monomial has zero Hurwitz-zeta regularized
sum because $\zeta(-2j,1/2)=0$.  Replacing $x_n$ by $n+1$ without
recomputing the asymptotic polynomial would introduce nonzero Bernoulli
counterterms and invalidate \eqref{eq:ordinary-reconstruction}.
Likewise, omitting the constant term of $Q_{b,k}$ is not harmless:
$\zeta(0,1/2)=0$ is exactly what permits that constant to be subtracted
without changing the regularized value.

\subsection{Novelty and literature scope}

The inspected project already contains extensive work on harmonic-tail
reconstruction, Stieltjes correlations, and multiple polylogarithms.
This report claims the following relative contributions only:

\begin{itemize}
\item the exact completion of the $p=10$ fixed-order family;
\item the explicit all-coordinate ordinary identities that follow from
that completion;
\item the two-tail generator \eqref{eq:two-tail-generator} in the
normalization used here; and
\item the pole-subtracted kernel \eqref{eq:K-kernel} and its coefficient
formula for integer-shift Stieltjes products.
\end{itemize}

A targeted literature search found the standard multiple-polylogarithm,
Hurwitz-zeta, Bernoulli, and Stieltjes ingredients, but not these exact
combined formulas in this normalization.  This is not a comprehensive
priority search, and no global first-discovery claim is made.  The
article remains an unrefereed research draft.

\subsection{Recommended manuscript placement}

The determinant theorem belongs immediately after the current
fixed-order and eventual-rank material.  The ordinary reconstruction
and polylogarithmic generator can follow the harmonic-tail chapter as a
new all-coordinate consequence.  The Stieltjes kernel is logically
independent of the order-ten certificate and may instead be placed in
the shifted-Stieltjes chapter, with a cross-reference to the tail
generator from which it is obtained by singular-order regularization.
The supplied integration fragment and claim ledger use namespaced labels
to avoid collisions during that editorial split.
